\documentclass[11pt]{amsart}
\usepackage{calc,amssymb,amsmath,ulem,stmaryrd,fullpage}
\usepackage{enumerate}
\usepackage{alltt}
\RequirePackage[dvipsnames,usenames]{color}
\let\mffont=\sf
\normalem
%\input{kmacros3.sty}
\input{xy}
\xyoption{all}
\usepackage{hyperref}
\usepackage{graphicx}
\usepackage[all,cmtip]{xy}
\usepackage{verbatim}
\usepackage[left=0.95in,top=0.95in,right=0.95in,bottom=0.95in]{geometry}
\newcommand{\tauCohomology}{T}
\newcommand{\FCohomology}{S}


\theoremstyle{theorem}
%\newtheorem*{mainthm*}{Main Theorem}

\newcommand{\note}[1]{\marginpar{\sffamily\tiny #1}}

\def\todo#1{\textcolor{Mahogany}%
{\sffamily\footnotesize\newline{\color{Mahogany}\fbox{\parbox{\textwidth-15pt}{#1}}}\newline}}

\renewcommand{\O}{\mathcal O}
\newcommand{\ie}{{\itshape i.e.} }
\newcommand{\roundup}[1]{\lceil #1 \rceil}
\newcommand{\rounddown}[1]{\lfloor #1 \rfloor}
\renewcommand{\to}[1][]{\xrightarrow{\ #1\ }}
\newcommand{\onto}[1][]{\protect{\xrightarrow{\ #1\ }\hspace{-0.8em}\rightarrow}}
\newcommand{\into}[1][]{\lhook \joinrel \xrightarrow{\ #1\ }}
%\newcommand{DBDual}[1]{{\underline{\omega}_{#1}}}
\newcommand{\DBDual}[1]{{{\uuline \omega}{}^{\mydot}_{#1}}}
\newcommand{\Tt}{{\mathfrak{T}}}
%\newcommand{\bn}{{\mathfrak{n}} }
\newcommand{\sol}[1]{ \vskip 6pt{\color{blue} {\bf Solution:} #1} \vskip 6pt}

\newcommand{\bR}{{\mathbb{R}}}
\newcommand{\va}{{\vec{a}}}
\newcommand{\vb}{{\vec{b}}}
\newcommand{\vzero}{{\vec{0}}}
\newcommand{\vi}{{\vec{i}}}
\newcommand{\vj}{{\vec{j}}}
\newcommand{\vk}{{\vec{k}}}
\newcommand{\vh}{{\vec{h}}}
\newcommand{\vr}{{\vec{r}}}
\newcommand{\vu}{{\vec{u}}}
\newcommand{\vv}{{\vec{v}}}
\newcommand{\vw}{{\vec{w}}}
\newcommand{\vx}{{\vec{x}}}
\newcommand{\vy}{{\vec{y}}}
\newcommand{\vz}{{\vec{z}}}
\newcommand{\vT}{{\vec{T}}}
\newcommand{\vN}{{\vec{N}}}
\newcommand{\vF}{{\vec{F}}}
\newcommand{\vG}{{\vec{G}}}
\newcommand{\bF}{\mathbb{F}}
\newcommand{\bQ}{\mathbb{Q}}
\newcommand{\bZ}{\mathbb{Z}}
\newcommand{\bC}{\mathbb{C}}
\newcommand{\Gal}{\textnormal{Gal}}
\newcommand{\Aut}{\textnormal{Aut}}
\newcommand{\Emb}{\textnormal{Emb}}
\newcommand{\Hom}{\textnormal{Hom}}


\begin{document}
\title{WS \#4 -- Math 6310\\Fall 2017}
\author{Due: Monday, December 4th}
\maketitle

You are encouraged to work in groups on this.  Only one write up needs to be turned in per group.  We learn about $\Hom$ in this worksheet.  Today, \emph{ALL} rings are commutative.
\vskip 6pt
\noindent
{\bf 1.}  Show that the functor $\Hom_R(\bullet, N)$ is left exact and contravariant.  In other words, if
\[
0 \to A \xrightarrow{f} B \xrightarrow{g} C \to 0
\]
is an exact sequence of $R$-modules, then
\[
0 \to \Hom_R(C, N) \xrightarrow{g'} \Hom_R(B, N) \xrightarrow{f'} \Hom_R(A, N)
\]
is exact.  (Did you use the fact that $A \to B$ was injective?)
\newpage
\noindent
{\bf 2.}  We prove a converse to {\bf 1.}   Suppose that $A \to B \to C$ is a sequence of maps between $R$-modules and also that $0 \to \Hom_R(C, N) \xrightarrow{g'} \Hom_R(B, N) \xrightarrow{f'} \Hom_R(A, N)$ is exact for \emph{every} $R$-module $N$.  Prove that $A \to B \to C \to 0$ is exact.
\vskip 3pt
\emph{Hint:} The trick is clever choices of $N$ to get different parts of the exactness.
\vskip 250pt
\noindent
{\bf 3.}  Show that the functor $\Hom_R(M, \bullet)$ is right exact and covariant.  In other words if
\[
0 \to A \xrightarrow{f} B \xrightarrow{g} C \to 0
\]
is an exact sequence of $R$-modules, then show that
\[
0 \to \Hom_R(M, A) \xrightarrow{f''} \Hom_R(M, B) \xrightarrow{g''} \Hom_R(M, C)
\]
is also exact.  (Did you use the fact that $B \to C$ was surjective?)
\newpage
\noindent
{\bf 4.}  Now we prove a converse to ${\bf 3.}$   Suppose that $A \to B \to C$ is a sequence of maps between $R$-modules and also that $0 \to \Hom_R(M, A) \xrightarrow{f''} \Hom_R(M, B) \xrightarrow{g''} \Hom_R(M, C)$ is exact for \emph{every} $R$-module $M$.  Prove that $0 \to A \to B \to C$ is exact.
\vskip 3pt
\emph{Hint:} Clever choices of $M$ are the order of the day.
\vskip 250pt
{\bf 5.}  A short exact sequence $0 \to A \xrightarrow{\alpha} B \xrightarrow{\beta} C \to 0$ is called \emph{split exact} if there is a map $\gamma : C \to B$ such that $\beta \circ \gamma : C \to B \to C$ is the identity.  Show that in that case $B \cong A \oplus C$ and $\alpha$ is identified with $a \mapsto (a, 0)$ and $\beta$ is identified with $(a,c) \mapsto c$.
\vskip 3pt
\emph{Fact:} In this case, there is also a map $\delta : B \to A$ so that $\delta \circ \alpha : A \to B \to A$ is the identity, which is also equivalent to showing that the short exact sequence is split exact.
\vskip 250pt
A $R$-module $P$ is called \emph{projective} if for every \emph{surjective} map of $R$-modules $f : A \to B$ and every $R$-module map $g : P \to B$, there exists a map $h : P \to A$ such that the following diagram commutes:
    \[
    \xymatrix{
    & P \ar@{.>}[dl]_{h} \ar[d]^g \\
    A \ar[r]_f & B.
    }
    \]
\noindent
{\bf 6.}    Show that a free module $P = R^{\oplus (...)}$ is projective.  Further show that if $P$ is projective and $P \cong A \oplus B$, then $A$ is also projective.
\vskip 200pt
\noindent
{\bf 7.}  Suppose that $0 \to A \xrightarrow{f} B \xrightarrow{g} P \to 0$ is exact and $P$ is projective, prove that
    \[
    0 \to \Hom_R(P, M) \to \Hom_R(B, M) \to \Hom_R(A, M) \to 0
    \]
    is exact for every $R$-module $M$.
    \vskip 3pt
    \emph{Hint: } Show that if $P$ is projective, then the exact sequence is split exact.
\end{document}

